% !TeX program = xelatex
% ============================================================
%  Causality: Models, Reasoning, and Inference (2nd ed.)
%  中文译稿 —— 共享导言区
% ============================================================
\usepackage{amsmath,amssymb,amsthm,mathrsfs,bm}
\usepackage{mathtools}
\usepackage{tikz}
\usetikzlibrary{arrows.meta,positioning,calc,fit,backgrounds,shapes.geometric,decorations.pathreplacing}
\usepackage{graphicx}
\usepackage{booktabs,array,multirow,longtable}
\usepackage{enumitem}
\usepackage{xcolor}
\usepackage[most]{tcolorbox}
\usepackage{caption}
\usepackage{ragged2e}
\usepackage{fancyhdr}
\usepackage{emptypage}

% ---------- 颜色 ----------
\definecolor{inkblue}{RGB}{28,62,110}
\definecolor{boxbg}{RGB}{246,246,242}
\definecolor{boxline}{RGB}{150,150,140}
\definecolor{ptbg}{RGB}{240,244,250}

% ---------- 页面 ----------
\usepackage[a4paper,top=2.6cm,bottom=2.6cm,left=2.6cm,right=2.6cm,headsep=0.7cm,footskip=1.0cm]{geometry}
\linespread{1.32}
\setlength{\parskip}{0.25em}

% ---------- 标题格式 ----------
\ctexset{
  chapter = {
    format     = \centering\heiti\zihao{2},
    name       = {第,章},
    number     = \chinese{chapter},
    beforeskip = 1.2cm,
    afterskip  = 1.2cm,
  },
  section = {
    format = \heiti\zihao{3},
    beforeskip = 1.1ex plus .3ex minus .2ex,
    afterskip  = .8ex plus .2ex,
  },
  subsection = {
    format = \heiti\zihao{4},
    beforeskip = 1.0ex plus .3ex minus .2ex,
    afterskip  = .6ex plus .2ex,
  },
  subsubsection = {
    format = \heiti\zihao{-4},
  },
  figurename = 图,
  tablename  = 表,
}

\captionsetup{font=small,labelsep=quad}
\captionsetup[figure]{justification=centering}

% ---------- 页眉页脚 ----------
\pagestyle{fancy}
\fancyhf{}
\fancyhead[LE]{\small\kaishu\leftmark}
\fancyhead[RO]{\small\kaishu\rightmark}
\fancyfoot[C]{\small\thepage}
\renewcommand{\headrulewidth}{0.4pt}
\fancypagestyle{plain}{\fancyhf{}\fancyfoot[C]{\small\thepage}\renewcommand{\headrulewidth}{0pt}}

% 页眉不做全大写（默认 \MakeUppercase 会把“do 算子”变成“DO 算子”）
\makeatletter
\renewcommand{\chaptermark}[1]{\markboth{\CTEXthechapter\quad #1}{}}
\renewcommand{\sectionmark}[1]{\markright{\thesection\quad #1}}
\makeatother

% ---------- 定理类环境（与原书一致：按节共享一个计数器）----------
\newtheorem{thmshared}{}[section]
\theoremstyle{definition}
\newtheorem{definition}[thmshared]{定义}
\newtheorem{example}[thmshared]{例}
\newtheorem{assumption}[thmshared]{假设}
\theoremstyle{plain}
\newtheorem{theorem}[thmshared]{定理}
\newtheorem{lemma}[thmshared]{引理}
\newtheorem{corollary}[thmshared]{推论}
\newtheorem{proposition}[thmshared]{命题}
\newtheorem{conjecture}[thmshared]{猜想}
\theoremstyle{remark}
\newtheorem{remark}[thmshared]{注}
\newenvironment{remarkenv}{\par\addvspace{0.6em}\noindent\heiti 注\quad\kaishu\ignorespaces}{\par\addvspace{0.4em}}

\renewcommand{\proofname}{\heiti 证明}
\renewcommand\qedsymbol{$\blacksquare$}

% ---------- 自定义盒子 ----------
% 小结（章内小结，替代原书 Summary 段）
\newtcolorbox{summarybox}[1][小结]{
  enhanced, breakable, colback=boxbg, colframe=boxline,
  boxrule=0.5pt, left=8pt, right=8pt, top=6pt, bottom=6pt,
  fonttitle=\heiti\small, title={#1},
  attach boxed title to top left={xshift=6pt,yshift*=-\tcboxedtitleheight/2},
  boxed title style={colback=boxbg,colframe=boxline,boxrule=0.5pt},
  before skip=10pt, after skip=10pt,
}
% 要点（对应原书的 fine point 式提示）
\newtcolorbox{keypoint}[1][要点]{
  enhanced, breakable, colback=white, colframe=inkblue,
  boxrule=0.6pt, leftrule=2.4pt, left=8pt, right=8pt, top=5pt, bottom=5pt,
  fonttitle=\heiti\small\color{inkblue}, title={#1},
  attach title to upper={\ \ },
  before skip=8pt, after skip=8pt,
}
% 技术注记
\newtcolorbox{techpoint}[1][技术注记]{
  enhanced, breakable, colback=ptbg, colframe=inkblue!45,
  boxrule=0.4pt, left=8pt, right=8pt, top=5pt, bottom=5pt,
  fonttitle=\heiti\small\color{inkblue}, title={#1},
  attach title to upper={\ \ },
  before skip=8pt, after skip=8pt,
}

% ---------- 数学算子 ----------
\DeclareMathOperator{\doop}{do}
\newcommand{\dox}[1]{\mathrm{do}(#1)}
\def\doxp(#1){\mathrm{do}(#1)}
\newcommand{\PA}{\mathrm{PA}}
\newcommand{\pa}{\mathrm{pa}}
\newcommand{\ND}{\mathrm{ND}}
\newcommand{\pr}{\Pr}
\newcommand{\given}{\,|\,}
\newcommand{\indep}{\perp\!\!\!\perp}
\newcommand{\nindep}{\not\perp\!\!\!\perp}
\newcommand{\deq}{\stackrel{\mathrm{def}}{=}}
\newcommand{\X}{\mathsf{X}}

% ---------- TikZ 统一风格 ----------
\tikzset{
  causal/.style={>={Stealth[length=2.1mm,width=1.5mm]},semithick},
  vtx/.style={circle,draw,inner sep=1pt,minimum size=6.2mm,font=\small},
  vtxb/.style={circle,draw,very thick,inner sep=1pt,minimum size=6.2mm,font=\small},
  vtxs/.style={circle,draw,inner sep=1pt,minimum size=5.4mm,font=\footnotesize},
  lat/.style={circle,draw,dashed,inner sep=1pt,minimum size=6.2mm,font=\small},
  ed/.style={->,semithick},
  edb/.style={->,very thick},
  edd/.style={->,semithick,dashed},
  bi/.style={<->,semithick},
  bid/.style={<->,semithick,dashed},
}

% ---------- 杂项 ----------
\newcommand{\en}[1]{\textup{#1}}
\renewcommand{\arraystretch}{1.15}
\setlist[itemize]{leftmargin=1.6em,itemsep=0.15em,topsep=0.3em}
\setlist[enumerate]{leftmargin=1.9em,itemsep=0.15em,topsep=0.3em}

% ---------- 超链接（最后加载）----------
\usepackage{url}
\usepackage[bookmarks=true,bookmarksnumbered=true,colorlinks=true,
            linkcolor=inkblue,citecolor=inkblue,urlcolor=inkblue]{hyperref}

% ---------- 断行宽容度（减少行溢出）----------
\setlength{\emergencystretch}{1.2em}
\tolerance=2000
